\abstract{%
Intrinsic-dimension estimates and active-unit diagnostics address different
aspects of variational-autoencoder bottleneck selection. We evaluate
search policies on four image datasets, five training seeds and five
regularisation weights, using a validation quality target and common
recorded candidate costs. A start-width ablation holds the reliability
checks, fallback and reference-bank charges fixed. The ID-derived and
midpoint starts return identical widths in all 100 conditions; where ID
is accepted, the former requires 8.84 rather than 5.84 candidates on average.
Two errors of an ungated midpoint policy occur in conditions where the
ID policy instead uses ascending fallback. Across 10,000 retrospective
coefficient and threshold settings, 155 local outputs miss the smallest
qualifying width. Additional three-seed experiments on each dataset
evaluate ARM-based pruning and exact-Jacobian relevance readouts.
Native reconstruction, dimension transfer and constraint satisfaction
give distinct outcomes; dSprites remains problematic for the local
ARM implementation. A separate one-seed MNIST experiment compares
Monte Carlo ELBO capacity selection. Active-unit prediction does not
establish additional benefit on its near-ceiling task. These results
identify how search rules, implementation choices and evaluation
definitions affect measured quality and cost, without establishing
universal dimension-selection guarantees.}
\keyword{variational autoencoder; bottleneck dimension; intrinsic dimension;
active units; model selection; reconstruction quality; reproducibility}
